\par\Needspace{6\baselineskip}\medskip
\begin{center}\hypertarget{native-169-article122}{}\pdfbookmark[2]{122.}{article-122}122.\end{center}
Die Werte, die man in der zweiten und dritten Hypothese für die Unbekannten ansetzt, sind in gewissem Maße beliebig. Sie dürfen sich nur nicht zu weit von denen der ersten Hypothese entfernen, und man muss darauf achten, dass sich das Verhältnis $(a''-a):(b''-b)$ dem Verhältnis $(a'-a):(b'-b)$ nicht zu sehr nähert. Meist setzt man daher $a'=a$, $b''=b$. Dies bringt einen doppelten Vorteil. Die Formeln für $\xi,\eta$ werden noch etwas einfacher, und ein Teil der ersten Rechnung bleibt bei der zweiten Hypothese, ein anderer bei der dritten derselbe.

Es gibt jedoch einen Fall, in dem andere Erwägungen dazu raten, von diesem Vorgehen abzuweichen. Nehmen wir an, $X$ habe die Form $X'-x$ und $Y$ die Form $Y'-y$, und die Funktionen $X',Y'$ seien nach der Natur der Aufgabe so beschaffen, dass sie von mäßigen Fehlern in den Werten von $x,y$ nur sehr wenig beeinflusst werden. Das heißt, $\left(\dfrac{\mathrm{d}X'}{\mathrm{d}x}\right)\mathpunct{\raisebox{.45ex}{,}}$ $\left(\dfrac{\mathrm{d}X'}{\mathrm{d}y}\right)\mathpunct{\raisebox{.45ex}{,}}$ $\left(\dfrac{\mathrm{d}Y'}{\mathrm{d}x}\right)\mathpunct{\raisebox{.45ex}{,}}$ $\left(\dfrac{\mathrm{d}Y'}{\mathrm{d}y}\right)$ sind sehr kleine Größen. Dann lassen sich die Unterschiede \SourceLeaf{170}zwischen den Werten dieser Funktionen, die dem Wertepaar $x=\xi$, $y=\eta$ entsprechen, und denjenigen, die sich aus $x=a$, $y=b$ ergeben, offenbar gewissermaßen einer höheren Ordnung zuweisen als die Unterschiede $\xi-a$, $\eta-b$. Die ersteren Werte sind $X'=\xi$, $Y'=\eta$, die letzteren $X'=a+A$, $Y'=b+B$. Daraus folgt, dass $a+A$, $b+B$ weit genauere Werte von $x,y$ sind als $a,b$. Legt man diese Werte der zweiten Hypothese zugrunde, so sind die Gleichungen $X=0$, $Y=0$ sehr oft schon so genau erfüllt, dass man nicht weiterzugehen braucht. Andernfalls bildet man aus der zweiten Hypothese auf dieselbe Weise die dritte, indem man $a''=a'+A'=a+A+A'$, $b''=b'+B'=b+B+B'$ setzt. Erweist sich diese noch als nicht genau genug, so gewinnt man daraus schließlich nach Artikel 120 eine vierte Hypothese.
